\chapter{精确密度模型：归一化流与自回归分解}
\section*{概述}
\addcontentsline{toc}{section}{概述}

前面几章已经展示了进入现代生成建模的多条主要路径。
变分自编码器引入了潜变量建模与近似推断；去噪扩散概率模型把视角转向随机加噪与去噪过程；连续时间与得分章节又进一步通过密度演化、随机微分方程和学习得到的得分场发展了这一思想。

本章转向另一种范式：\emph{精确密度建模}。
这里的目标是构造这样一类生成模型：对于一个观测数据点，其概率密度在模型设计上就可以直接计算。
这改变了建模哲学。
我们不再依赖难以处理的潜变量边缘似然，也不再学习去反转随机扰动过程，而是主动施加某种数学结构，使似然计算从一开始就保持可处理。

第一条路线是\textbf{归一化流}（normalising flows）。
归一化流从一个简单基础分布（例如标准高斯）出发，依次应用一系列可微、可逆变换。
变量变换公式通过 Jacobian 行列式刻画局部拉伸或压缩，从而精确告诉我们密度在每一次变换下如何变化。
只要这些变换可逆，且其 Jacobian 行列式可处理，那么似然计算和采样都是模型中的精确操作。

第二条路线是\textbf{自回归模型}。
它不通过可逆映射对基础分布进行变换，而是直接利用概率链式法则构造联合密度。
一个高维数据点被拆解成有序分量序列，每个分量都以此前分量为条件进行建模。
这种方式无需潜变量、变分下界或 Jacobian 行列式，也能得到精确似然，但采样通常必须顺序执行。

这两种方法在现代生成建模中仍然十分重要，因为它们展示了精确似然建模能够达到怎样的程度。
归一化流为可逆密度估计、不确定性建模以及概率分布的可处理变换提供了清晰框架。
与此同时，自回归模型构成许多有影响力的序列、音频、语言以及像素级图像生成系统的基础。
即使较新的系统会组合多种建模思想，本章涉及的原则——变量变换、可处理 Jacobian、条件分解以及有序生成——仍然会作为核心构件不断出现。

\section*{概念导图}
\addcontentsline{toc}{section}{概念导图}

本章把精确密度建模作为与此前潜变量、扩散和得分方法相对照的另一种范式展开。
核心问题是：灵活的高维生成模型如何同时保留精确似然计算能力？

\begin{itemize}
    \item[\(\rightarrow\)] \textbf{首先通过变量变换公式引入归一化流。}
    归一化流利用可微、可逆变换把简单基础变量 $z$ 映射为数据变量 $x=f(z)$。
    变量变换公式通过 Jacobian 行列式跟踪相应的密度变化。

    \item[\(\rightarrow\)] \textbf{随后把平面流作为第一种非线性流层来研究。}
    平面流通过秩一残差扰动引入非线性，在保持 Jacobian 行列式解析可处理的同时，也展示了可逆性约束的重要性。

    \item[\(\rightarrow\)] \textbf{接着转向仿射耦合层这一实用流构造。}
    仿射耦合层只变换输入的一部分，而以另一部分为条件。
    这种设计使逆映射具有闭式形式，且对数行列式能够高效计算，因此适用于实际流模型。

    \item[\(\rightarrow\)] \textbf{然后把多个流层组合成完整生成模型。}
    堆叠许多可逆层后，简单局部形变可以组合成表达力很强的全局变换。
    似然计算把数据逆向映射回基础空间，而采样则把基础噪声正向映射到数据空间。

    \item[\(\rightarrow\)] \textbf{之后引入自回归分解，作为获得精确似然的第二条路线。}
    自回归模型利用链式法则
    $p(x)=\prod_{i=1}^d p(x_i\mid x_{<i})$
    构造精确似然，无需潜变量、变分下界或 Jacobian 行列式。

    \item[\(\rightarrow\)] \textbf{最后解释条件参数化、训练、采样与掩码。}
    共享模型可以参数化所有条件分布。
    训练时完整数据向量是已知的，因此似然项常可并行计算；采样时，各分量则必须顺序生成。
    掩码用来强制实现使自回归分解可处理的有序信息流。
\end{itemize}

本章的两条路线以及使其保持可处理性的结构原则总结如下图所示。

\begin{center}
\includegraphics[width=0.96\textwidth]{figures/ch07/chapter7_exact_density_at_a_glance.pdf}
\end{center}


\section{归一化流：通过可逆变换实现精确密度}

归一化流通过在简单参考分布与数据分布之间构造显式映射，同时始终保留精确似然计算能力，为生成建模提供了一条直接且完全概率化的路线~\cite{rezende2015variational}。
“归一化流”这一术语指的是这一更广泛的可逆密度模型家族，而不是某一种具体变换；平面流、仿射耦合层以及相关架构，都是设计这类变换的不同方式，其共同目标是同时保持表达力、可逆性与计算可处理性。

归一化流背后的核心数学工具是\emph{变量变换公式}。
该公式精确刻画概率密度在可微、可逆映射下如何变化。
归一化流不引入潜变量或辅助随机过程，而完全依赖确定性变换，并且这些变换对密度的影响可以精确计算。
通过串联许多这样的变换，一个简单基础分布——通常是标准高斯——可以逐渐被重塑成高度结构化的数据分布。

从几何角度看，每个变换都是对空间进行受控形变，通过局部拉伸或压缩体积来改变概率密度，而这种变化由 Jacobian 行列式显式计入。
只要每个变换可逆且其 Jacobian 行列式可处理，所得模型就同时支持精确似然计算和精确采样。
这种表达力与数学透明性的结合，使归一化流区别于其他生成建模范式，也使其成为研究精确密度模型的自然起点。

\subsection{重新审视变量变换公式}

虽然变量变换公式已经在第~5 章从更一般的确定性密度输运角度介绍过，但这里仍值得重新审视，因为它正是归一化流的建模原则。

假设 $X$ 有已知密度 $p_X(x)$，并应用一个可微、可逆函数
$Y=f(X)$。
我们的目标是确定 $Y$ 的密度 $p_Y(y)$。
直观上，$Y$ 的密度需要同时考虑两件事。

第一，它必须考虑\emph{概率质量从哪里来}。
由于 $y=f(x)$，观测到 $y$ 的可能性取决于其对应原像
$x=f^{-1}(y)$
在原分布下有多大概率。

第二，它必须考虑\emph{变换如何拉伸或压缩空间}。
如果一个变换让点彼此离得更远，密度应当降低；如果让点压缩得更紧，密度应当提高。
密度始终是单位长度、面积或体积上的概率，因此局部间距至关重要。

在一维中，这两个思想结合为变量变换公式
\[
\boxed{
p_Y(y)
=
p_X(f^{-1}(y))
\left|
\frac{d f^{-1}(y)}{dy}
\right|.
}
\]
第一项告诉我们哪个 $x$ 生成了 $y$；第二项是缩放因子，用来补偿 $x$ 附近的邻域映射到 $y$ 后被拉伸或压缩的程度。

\paragraph{一个简单几何例子。}

设
$X\sim \mathrm{Uniform}(0,1),
\qquad
Y=f(X)=2X$。

这个变换把所有点扩大两倍，把区间 $[0,1]$ 拉伸成 $[0,2]$。
其逆映射为
\[
x=f^{-1}(y)=\frac{y}{2},
\]
所以
\[
\frac{d f^{-1}(y)}{dy}
=
\frac{1}{2}.
\]
代入公式得到
\[
p_Y(y)
=
p_X\!\left(\frac{y}{2}\right)\frac{1}{2}.
\]
在新区间 $[0,2]$ 内，原密度满足 $p_X(y/2)=1$，因此
\[
p_Y(y)=\frac{1}{2},
\qquad 0\leq y\leq 2.
\]
这与直觉完全一致：同样的总概率现在填充一个宽度变为两倍的区间，所以密度减半。
这个简单逻辑可以直接推广到高维。

\paragraph{高维情形：从长度拉伸到体积拉伸。}

在多维中，普通导数被体积缩放的行列式取代。
对于简单线性映射
\[
f(x)=Ax,
\qquad
A=
\begin{bmatrix}
2 & 0\\
0 & 3
\end{bmatrix},
\]
局部面积会乘以 $|\det A|=6$，因此
\[
p_Y(y)=p_X(A^{-1}y)\frac{1}{|\det A|}.
\]
同样的概率质量占据了更大的局部体积，所以密度按照其倒数因子下降。

\paragraph{Jacobian：观察局部线性变化的透镜。}

为了推广到非线性变换，需要使用\emph{Jacobian 矩阵}。
对于可微函数
\[
f:\mathbb{R}^d\to\mathbb{R}^d,
\]
Jacobian 定义为
\[
[J_f(x)]_{ij}
=
\frac{\partial f_i(x)}{\partial x_j}.
\]
可以把 Jacobian 看作观察局部线性行为的一面透镜。
在点 $x$ 附近，一个很小的位移 $\delta x$ 近似按照
\[
f(x+\delta x)
\approx
f(x)+J_f(x)\delta x
\]
变换。
因此，$J_f(x)$ 告诉我们 $x$ 附近一个无穷小邻域映射到 $y=f(x)$ 后如何发生形变；其行列式则量化对应局部体积膨胀或收缩了多少。

对于一般可微可逆变换 $y=f(x)$，多元变量变换公式为
\[
\boxed{
p_Y(y)
=
p_X(f^{-1}(y))
\left|
\det J_{f^{-1}}(y)
\right|
=
p_X(x)
\frac{1}{|\det J_f(x)|},
\qquad y=f(x).
}
\]

这就是归一化流的数学核心：
\emph{只要每个变换可逆且其 Jacobian 行列式可处理，即使整体映射具有很强表达力，我们仍然能够精确写出变换后变量的密度。}
图~\ref{fig:flow-change-variable-geometry} 展示了这一公式的几何含义：固定的概率质量落在局部被拉伸或压缩的区域中，密度按体积变化因子的倒数进行调整。

\begin{figure}[H]
\centering
\includegraphics[width=0.88\textwidth]{figures/ch07/figure7_1_change_of_variables_geometry.pdf}
\caption{\textbf{变量变换通过局部体积变化修正概率密度。}
同一份概率质量经过可逆变换后会占据不同的局部面积。在线性例子中，行列式给出面积缩放比例；对于非线性流，Jacobian 则逐点给出相应的局部缩放。这就是对数行列式校正项的几何来源。}
\label{fig:flow-change-variable-geometry}
\end{figure}


\begin{tcolorbox}[
  colback=blue!4,
  colframe=blue!50!black,
  title={为什么标准归一化流保持维度不变},
  boxrule=0.6pt,
  arc=2mm,
  left=2mm,
  right=2mm,
  top=1mm,
  bottom=1mm
]
标准归一化流采用形式为
\[
f:\mathbb{R}^d\to\mathbb{R}^d
\]
的可微可逆变换。
相同的输入输出维度是标准变量变换公式所要求的：Jacobian 必须是方阵，$\det J_f(x)$ 才有定义，并能够度量局部体积变化。
如果映射把 $\mathbb{R}^d$ 映到更低维空间，信息会被压塌，不同输入可能映射到同一个输出；如果把 $\mathbb{R}^d$ 映到更高维空间，像会落在低维流形上，在没有额外结构时并不能定义普通的满维密度。
因此，标准归一化流保持维度不变。
若失去这种满维可逆性，就无法使用变量变换公式精确追踪概率质量。
\end{tcolorbox}

\subsection{平面流}

上述例子是有意选择的简单情形，它们展示了变量变换公式如何跟踪线性拉伸与体积缩放下的密度。
然而，线性变换的表达能力有限：它们可以拉伸、旋转或剪切空间，却无法以丰富的非线性方式弯曲或扭曲空间。
为了逼近复杂分布，我们需要能够引入受控非线性，同时仍保持归一化流两项核心要求的变换：可逆性与可处理的 Jacobian 行列式。

平面流（planar flow）是这种权衡的经典例子，最初作为 Rezende 与 Mohamed~\cite{rezende2015variational} 的归一化流框架的一部分提出。
一种基础非线性流层定义为
\[
f(x)
=
x+u\,h(w^\top x+b),
\]
其中 $x\in\mathbb{R}^d$，$u,w\in\mathbb{R}^d$，$b\in\mathbb{R}$，且
\[
h:\mathbb{R}\to\mathbb{R}
\]
是光滑激活函数，例如 $\tanh$ 或 \texttt{sigmoid}。
这里 $x$ 表示当前单个流层的输入；在堆叠流中，它通常是某个中间变量。

项
\[
w^\top x+b
\]
对输入进行仿射投影，把高维向量 $x$ 压缩为一个标量。
单独使用这个投影只会产生线性量。
为了引入曲率和更丰富的行为，我们把该标量送入非线性激活函数 $h(\cdot)$。
$\tanh$ 与 \texttt{sigmoid} 之类的函数会把输入压缩到有界范围，因此 $h$ 的输出可以看成一个软门控值，用来控制变换应该多强地修改输入。

把这个门控标量与向量 $u$ 相乘，会产生一个具有方向的缩放扰动。
形变始终沿 $u$ 的方向发生，但其幅度以非线性方式依赖投影 $w^\top x+b$。
把这一扰动重新加回 $x$，既保证输出仍然位于完整的 $d$ 维空间中，又引入了灵活的非线性修正。

\begin{tcolorbox}[
  colback=blue!4,
  colframe=blue!50!black,
  title={残差结构及其动机},
  boxrule=0.6pt,
  arc=2mm,
  left=2mm,
  right=2mm,
  top=1mm,
  bottom=1mm
]
平面流具有深度学习中残差连接的基本精神：
\[
f(x)=x+\text{学习得到的扰动}.
\]
它保留恒等映射，同时叠加受控形变。

如果只使用非线性项
\[
x\mapsto u\,h(w^\top x+b),
\]
由于 $h(\cdot)$ 输出标量，结果会塌缩到 $u$ 所在的一维方向，从而破坏归一化流所需的满维结构。

把扰动加回 $x$ 后，平面流能够保持整个 $d$ 维空间可达，同时引入有针对性的非线性形变。
\end{tcolorbox}

\paragraph{Jacobian 推导。}
为了计算平面流作用下密度如何变换，我们推导
\[
f(x)=x+u\,h(w^\top x+b)
\]
的 Jacobian。为简化记号，令
\[
a=w^\top x+b,
\qquad
f(x)=x+u\,h(a).
\]
对 $x$ 求导并使用链式法则，得到
\[
\frac{\partial f(x)}{\partial x}
=
I+u\,h'(a)\frac{\partial a}{\partial x}.
\]
由于 $\frac{\partial a}{\partial x}=w^\top$，Jacobian 可写成紧凑形式
\[
J_f(x)
=
I+u\,w^\top h'(a).
\]
等价地，
\[
J_f(x)
=
I+u\psi(x)^\top,
\qquad
\psi(x):=h'(w^\top x+b)w.
\]

这一结构非常便利。Jacobian 是单位矩阵加上一个秩一矩阵扰动。因此，无须计算通常需要 $O(d^3)$ 代价的完整 $d\times d$ 行列式，行列式可以化为单个标量表达式。利用矩阵行列式引理
\[
\det(I+uv^\top)=1+v^\top u,
\]
可得
\[
\det J_f(x)
=
1+\psi(x)^\top u.
\]
代入 $\psi(x)=h'(w^\top x+b)w$，得到
\[
\boxed{
\det J_f(x)
=
1+h'(w^\top x+b)\,w^\top u.
}
\]
因此，平面流对对数行列式的贡献为
\[
\log|\det J_f(x)|
=
\log\left|
1+h'(w^\top x+b)\,w^\top u
\right|.
\]
这一闭式表达式使平面流能够用于基于似然的可处理训练。

\begin{tcolorbox}[
  colback=blue!4,
  colframe=blue!50!black,
  title={理解秩一更新},
  boxrule=0.6pt,
  arc=2mm,
  left=2mm,
  right=2mm,
  top=1mm,
  bottom=1mm
]
矩阵 $u w^\top$ 是 $\mathbb{R}^d$ 中两个向量的外积，而这样的矩阵秩为一。

要理解这一点，可将它作用于任意向量 $v$：
\[
u w^\top v
=
u\,(w^\top v).
\]
标量 $w^\top v$ 会随 $v$ 变化，但输出始终只是 $u$ 的某个标量倍数。因此，该矩阵只有一个独立的输出方向。

这也解释了为什么行列式如此简单：平面流并不试图同时对所有维度实施任意的非线性形变，而是引入一个受控的单方向形变，其体积变化可以解析计算。
\end{tcolorbox}

图~\ref{fig:planar-flow-rank-one} 展示了相应的几何图景：形变幅度随空间位置变化，但残差位移始终平行于单一方向 $u$。

\begin{figure}[H]
\centering
\includegraphics[width=0.92\textwidth]{figures/ch07/figure7_2_planar_flow_rank_one.pdf}
\caption{\textbf{平面流是非线性的，但其方向受到约束。}
位移大小通过标量非线性响应随位置变化，但所有残差位移始终平行于同一个向量 $u$。这正是正文所推导秩一结构的几何结果。}
\label{fig:planar-flow-rank-one}
\end{figure}


\paragraph{可逆性条件。}

除了 Jacobian 行列式必须可处理之外，流层还必须可逆。对于平面流
\[
f(x)=x+u\,h(w^\top x+b),
\]
其逆映射一般没有闭式表达式，因此通常通过参数约束来控制可逆性。

由前面的 Jacobian 推导，
\[
\det J_f(x)
=
1+h'(w^\top x+b)\,w^\top u.
\]
如果该行列式变为零，变换会在局部压缩体积，因而无法局部求逆。因此，一个常见的充分条件是让该行列式对所有 $x$ 都保持为正。

对于 $\tanh$ 这类满足
\[
0\leq h'(a)\leq 1
\]
的激活函数，只需要求
\[
w^\top u>-1.
\]
于是对任意 $a$ 都有
\[
1+h'(a)\,w^\top u>0,
\]
从而平面流可避免局部体积坍缩。实践中，可以通过对 $u$ 进行重参数化来强制满足这一条件。

因此，平面流很好地体现了一项核心设计权衡：通过简单的残差扰动引入非线性，同时使 Jacobian 行列式保持解析可处理，并通过一个标量条件控制可逆性。


\section{实践中的归一化流：耦合层与组合}

上一节将归一化流介绍为建立在可逆变换之上的精确密度模型。平面流给出了一个简单的非线性流层示例，其 Jacobian 行列式仍然易于处理，但也暴露出一个实际限制：其逆映射通常没有闭式形式。

现在，我们从这一教学性构造转向实用流架构背后的核心设计原则之一：\emph{耦合（coupling）}。耦合层每次只变换输入的一部分，并把剩余部分用作条件信号，由此保留精确可逆性。这种结构会产生三角 Jacobian，给出闭式逆映射，并允许高效堆叠多个层。

随后我们将说明如何把这些层组合成一个完整的生成模型。正是这种组合，把简单的局部变换转化为富有表达能力的密度模型，同时保留归一化流的两项定义性操作：精确似然计算，以及从学习到的模型中进行精确采样。

\subsection{仿射耦合层与三角 Jacobian}

仿射耦合层最早用于 NICE、RealNVP 等基于耦合的流模型~\cite{dinh2014nice,dinh2017density}，是构建具有可处理逆映射和对数行列式的高表达能力流模型时最实用的方法之一。

仿射耦合层首先把输入向量划分为两部分。设 $x\in\mathbb{R}^d$ 是某个流层的输入，并写为
\[
x=(x_1,x_2),
\]
其中 $x_1\in\mathbb{R}^{d_1}$、$x_2\in\mathbb{R}^{d_2}$，且 $d_1+d_2=d$。仿射耦合变换定义为
\[
\begin{aligned}
y_1 &= x_1,\\
y_2 &= x_2\odot \exp(s(x_1)) + t(x_1),
\end{aligned}
\]
其中 $s(\cdot)$ 与 $t(\cdot)$ 是函数，通常由神经网络表示，$\odot$ 表示逐元素乘法。

输入的第一部分保持不变。第二部分通过一个逐元素仿射映射进行变换，其缩放量和平移量由第一部分预测。因此，$x_1$ 充当上下文变量，而 $x_2$ 则依据这一上下文进行缩放和平移。

虽然该变换关于 $x_2$ 是仿射的，但函数 $s(x_1)$ 和 $t(x_1)$ 可以具有高度非线性。这使整体变换既有较强表达能力，又保持简单的代数结构。因此，该层有意采用非对称设计：牺牲单层内部所有维度的同时混合，换取精确求逆与高效行列式计算。

\paragraph{Jacobian 行列式。}

为了计算 Jacobian，写成
\[
y=
\begin{bmatrix}
y_1\\
y_2
\end{bmatrix}
=
\begin{bmatrix}
x_1\\
x_2\odot \exp(s(x_1))+t(x_1)
\end{bmatrix}.
\]
其 Jacobian 具有分块形式
\[
J
=
\frac{\partial y}{\partial x}
=
\begin{bmatrix}
\frac{\partial y_1}{\partial x_1}
&
\frac{\partial y_1}{\partial x_2}
\\[4pt]
\frac{\partial y_2}{\partial x_1}
&
\frac{\partial y_2}{\partial x_2}
\end{bmatrix}.
\]

由于 $y_1=x_1$，有
\[
\frac{\partial y_1}{\partial x_1}=I,
\qquad
\frac{\partial y_1}{\partial x_2}=0.
\]
左下角分块 $\frac{\partial y_2}{\partial x_1}$ 可能十分复杂，因为其中包含神经网络 $s(\cdot)$ 与 $t(\cdot)$ 的导数。然而关键在于，这一分块并不会影响行列式。

对于右下角分块，注意 $y_2$ 的每个分量为
\[
y_{2,i}
=
x_{2,i}\exp(s_i(x_1))+t_i(x_1).
\]
在固定 $x_1$ 的情况下，对 $x_{2,j}$ 求导可得
\[
\frac{\partial y_{2,i}}{\partial x_{2,j}}
=
\begin{cases}
\exp(s_i(x_1)), & i=j,\\
0, & i\neq j.
\end{cases}
\]
因此
\[
\frac{\partial y_2}{\partial x_2}
=
\operatorname{diag}\!\big(\exp(s(x_1))\big).
\]

把这些分块组合起来，得到
\[
J
=
\begin{bmatrix}
I & 0\\
* & \operatorname{diag}\!\big(\exp(s(x_1))\big)
\end{bmatrix}.
\]
这是一个分块下三角矩阵。对于任意三角矩阵或分块三角矩阵，其行列式等于对角分块行列式之积。因此
\[
\det J
=
\det(I)\,
\det\!\left(
\operatorname{diag}\!\big(\exp(s(x_1))\big)
\right).
\]
由于 $\det(I)=1$，而对角矩阵的行列式就是各对角元素之积，所以
\[
\det J
=
\prod_i \exp(s_i(x_1)).
\]
取对数得到特别简单的表达式
\[
\boxed{
\log|\det J|
=
\sum_i s_i(x_1).
}
\]

这正是仿射耦合层的主要可处理性来源。稠密且可能很复杂的分块 $\frac{\partial y_2}{\partial x_1}$ 无须参与行列式计算，全部体积变化都由缩放输出 $s(x_1)$ 控制。因此，即使层内部使用了非线性神经网络，对数行列式仍可以关于被变换维数在线性时间内计算。

三角分块结构是这里的关键计算简化：非对角分块可以很复杂，但不会进入行列式。这体现了精确密度建模中的一个反复出现的原则：以精心设计的方式组织非线性依赖，使高维概率计算仍能分解为简单项。

\paragraph{可逆性。}

带来三角 Jacobian 的同一种非对称结构，也同时给出了闭式逆映射。由前向变换
\[
\begin{aligned}
y_1 &= x_1,\\
y_2 &= x_2\odot \exp(s(x_1)) + t(x_1),
\end{aligned}
\]
可以立即恢复
\[
x_1=y_1.
\]
将其代入第二个方程，得到
\[
y_2
=
x_2\odot \exp(s(y_1)) + t(y_1).
\]
由于 $\exp(s(y_1))$ 逐分量严格为正，这个逐元素仿射变换可直接求逆：
\[
x_2
=
\bigl(y_2-t(y_1)\bigr)\odot \exp(-s(y_1)).
\]
因此，逆映射为
\[
\boxed{
\begin{aligned}
x_1 &= y_1,\\
x_2 &= \bigl(y_2-t(y_1)\bigr)\odot \exp(-s(y_1)).
\end{aligned}
}
\]

这里不需要求根或迭代求逆。逆变换与前向变换具有基本相同的计算代价。这使仿射耦合层尤其适合大规模模型，因为采样与似然计算往往需要反复应用许多流变换。图~\ref{fig:affine-coupling-structure} 总结了使这两种性质同时成立的非对称信息流：一部分变量原样通过，同时控制另一部分变量的变换。

\begin{figure}[H]
\centering
\begin{tikzpicture}[
  >=Latex,
  every node/.style={font=\small},
  state/.style={draw,rounded corners=2pt,align=center,inner sep=5pt,text width=2.35cm,minimum height=.82cm},
  note/.style={draw,rounded corners=2pt,align=center,inner sep=5pt,text width=4.45cm,minimum height=.82cm},
  flow/.style={->,thick}
]
\node[font=\normalsize\bfseries] at (6.0,2.75) {仿射耦合：单向依赖带来可处理性};

\node[state] (x1) at (.8,1.20) {$x_1$\\上下文};
\node[state] (x2) at (.8,-.55) {$x_2$\\被变换};
\node[state] (st) at (4.0,.30) {缩放与平移\\由 $x_1$ 决定};
\node[state] (y1) at (7.35,1.20) {$y_1$\\保持不变};
\node[state] (y2) at (7.35,-.55) {$y_2$\\被变换};

\draw[flow] (x1) -- (st);
\draw[flow] (x1) -- (y1);
\draw[flow] (st) -- (y2);
\draw[flow] (x2) -- (y2);

\node[note] (tri) at (3.0,-2.45) {单向依赖\\$\Rightarrow$ 分块三角 Jacobian};
\node[note] (inv) at (8.7,-2.45) {保持不变的分支\\$\Rightarrow$ 可在逆变换中直接恢复};
\end{tikzpicture}
\caption{\textbf{仿射耦合中的非对称信息流。}
一部分变量原样通过，同时决定另一部分变量如何变换。这种单向依赖同时带来了三角 Jacobian 和闭式逆映射。}
\label{fig:affine-coupling-structure}
\end{figure}


\subsection{堆叠流：似然计算与采样}

单个流层只能产生有限的空间形变。平面流沿一个学习到的方向对恒等映射施加扰动，而仿射耦合层每次只变换输入的一部分。因此，归一化流的表达能力来自\emph{组合}：许多简单、可逆的变换相互堆叠，使各自的小幅形变累积为丰富的全局映射。

设
\[
f_1,f_2,\dots,f_K
\]
为一系列可微且可逆的变换。从基础变量
\[
z_0\sim p_Z(z_0)
\]
出发，定义中间变量
\[
z_k=f_k(z_{k-1}),
\qquad k=1,\dots,K,
\]
以及最终输出
\[
x=z_K.
\]
等价地，
\[
x
=
f_K\circ f_{K-1}\circ \cdots \circ f_1(z_0).
\]
由于每一层都可逆，整个组合也可逆。因此我们能够沿两个方向移动：
\[
z_0
\xrightarrow{f_1}
z_1
\xrightarrow{f_2}
\cdots
\xrightarrow{f_K}
z_K=x,
\]
或者
\[
x=z_K
\xrightarrow{f_K^{-1}}
z_{K-1}
\xrightarrow{f_{K-1}^{-1}}
\cdots
\xrightarrow{f_1^{-1}}
z_0.
\]

这种双向结构给出了流模型的两个基本操作：似然计算与采样。

\paragraph{似然计算：从数据空间到基础空间。}

在基于似然的训练中，我们给定一个数据点 $x$。按照上述约定，$x$ 是最终变量 $z_K$。为了计算其密度，我们通过逆变换逐层向后映射：
\[
x=z_K
\xrightarrow{f_K^{-1}}
z_{K-1}
\xrightarrow{}
\cdots
\xrightarrow{f_1^{-1}}
z_0.
\]
如果模型拟合良好，恢复出来的基础空间变量 $z_0$ 应当在简单基础密度 $p_Z$ 下具有较高概率。

逐层应用变量变换公式，得到
\[
p_{Z_k}(z_k)
=
p_{Z_{k-1}}(z_{k-1})
\left|
\det J_{f_k}(z_{k-1})
\right|^{-1}.
\]
从 $k=1$ 到 $K$ 反复使用这一递推，并利用 $z_K=x$，得到
\[
\boxed{
p_X(x)
=
p_Z(z_0)
\prod_{k=1}^K
\left|
\det J_{f_k}(z_{k-1})
\right|^{-1}.
}
\]
取对数可得
\[
\boxed{
\log p_X(x)
=
\log p_Z(z_0)
-
\sum_{k=1}^K
\log
\left|
\det J_{f_k}(z_{k-1})
\right|.
}
\]

这一表达式具有直接的解释。基础密度项 $\log p_Z(z_0)$ 衡量经过逆变换的数据点在简单参考分布下有多合理；对数行列式项则校正整个变换序列所引入的全部体积变化。每一层都会向最终对数似然贡献一个可加的体积校正项。

\begin{tcolorbox}[
  colback=blue!3,
  colframe=blue!50!black,
  title={乘积形式是如何产生的},
  boxrule=0.6pt,
  arc=2mm,
  left=2mm,
  right=2mm,
  top=1mm,
  bottom=1mm
]
对于单个流层 $z_k=f_k(z_{k-1})$，变量变换公式给出
\[
p_{Z_k}(z_k)
=
p_{Z_{k-1}}(z_{k-1})
\left|\det J_{f_k}(z_{k-1})\right|^{-1}.
\]
反复应用这一关系，得到
\[
\begin{aligned}
p_{Z_2}(z_2)
&=
p_{Z_0}(z_0)
\left|\det J_{f_1}(z_0)\right|^{-1}
\left|\det J_{f_2}(z_1)\right|^{-1},\\
&\ \vdots\\
p_{Z_K}(z_K)
&=
p_{Z_0}(z_0)
\prod_{k=1}^K
\left|\det J_{f_k}(z_{k-1})\right|^{-1}.
\end{aligned}
\]
由于 $z_K=x$，这就给出了观测数据点的似然。
\end{tcolorbox}

\paragraph{采样：从基础空间到数据空间。}

采样沿相反方向使用同一组变换。首先从基础分布中抽取一个点，
\[
z_0\sim p_Z(z_0),
\]
然后依次应用学习到的变换：
\[
z_0
\xrightarrow{f_1}
z_1
\xrightarrow{f_2}
\cdots
\xrightarrow{f_K}
z_K=x.
\]
所得 $x$ 就是来自学习到的模型分布的一个样本。

相对于模型本身而言，这一过程是精确的：不需要 Markov 链、变分近似，也不需要数值积分。变换学习完成后，生成过程只是从基础分布采样，再应用一系列确定性的可逆映射。

这也进一步澄清了它与扩散模型的差别。扩散模型通过模拟从噪声到数据的学习到的反向随机过程来生成样本，而归一化流则应用一个学习到的确定性可逆变换。

\paragraph{组合、混合与表达能力。}

堆叠许多流层能够提升模型表达能力，但仅有组合还不够。层的排列还必须使信息能够在不同维度之间混合。如果每一层始终采用相同的划分，一些变量就会反复充当条件变量，却始终没有被直接变换。

为了解决这一问题，实用流模型通常会在非线性流层之间穿插简单的可逆混合操作。这些操作可以是置换
\[
z\mapsto Pz,
\]
也可以是更一般的可逆线性变换
\[
z\mapsto Az,
\]
其中 $P$ 是置换矩阵，$A$ 是具有可处理行列式的可逆矩阵。

这些混合操作的作用，是改变后续层中哪些坐标彼此交互。在一个耦合层中保持不变的变量，可以在下一层被移动到需要变换的子集中。从几何上看，每一次混合操作都改变了下一次非线性形变所处的坐标系。

因此，一个现代流模型通常由以下模式反复构成：
\[
\text{混合坐标}
\quad\longrightarrow\quad
\text{施加可处理的非线性形变}
\quad\longrightarrow\quad
\text{再次混合}.
\]
这种不断交替使局部、结构化的变换可以组合为灵活的全局映射，同时保留精确可逆性和可处理的似然计算。

RealNVP 和 Glow 等架构都建立在这一原则之上。它们在耦合变换形式、缩放参数化和坐标混合细节上有所不同，但底层数学结构相同：组合简单的可逆层，追踪每一层的体积变化，并通过精确最大似然训练整个模型。

\section{作为生成模型的自回归分解}

归一化流通过在简单基础分布与数据分布之间构造可逆映射，实现精确密度建模。其可处理性来自变量变换公式：只要变换可逆且 Jacobian 行列式易于计算，似然就可以被精确求得。

自回归模型则提供了第二条实现精确密度建模的路径。它不把基础随机变量变换为数据，而是直接把联合密度分解为一系列条件密度的乘积。此时，精确似然来自概率论的链式法则。这也改变了结构约束施加的位置：在归一化流中，结构施加在变换本身；在自回归模型中，结构施加在一个数据点各分量之间的依赖关系上。通过选定一个顺序，并规定每个分量只能依赖先前分量，复杂的联合密度就变成了若干较简单条件项的乘积。

因此，一个高维数据点可以被看作按顺序逐个揭示的一列分量。每个分量都在先前分量的条件下建模，由此得到的条件分布共同定义了一个有效联合分布，其似然可以精确计算。由此，自回归模型在不使用潜变量、变分下界或全局可逆变换的情况下，也构成了一套完整的生成建模框架。它付出的代价不同：必须显式规定顺序，并且采样过程变为顺序执行。

\subsection{作为精确密度模型的链式法则}

设
\[
x=(x_1,x_2,\dots,x_d)
\]
是一个 $d$ 维随机向量。概率链式法则可由恒等式
\[
p(a,b)=p(a\mid b)p(b)
\]
反复应用得到。我们可以从完整联合密度出发，先把最后一个变量分离出来：
\[
p(x_1,\dots,x_d)
=
p(x_d\mid x_1,\dots,x_{d-1})p(x_1,\dots,x_{d-1}).
\]
再对剩余联合项递归应用同一个恒等式，得到
\[
\begin{aligned}
p(x_1,\dots,x_d)
&=
p(x_d\mid x_1,\dots,x_{d-1})
p(x_{d-1}\mid x_1,\dots,x_{d-2})
\cdots
p(x_2\mid x_1)p(x_1).
\end{aligned}
\]
重新排列各个因子，即得到标准的自回归分解：
\[
\boxed{
p(x_1,\dots,x_d)
=
\prod_{i=1}^d p(x_i\mid x_{<i}),
}
\]
其中
\[
x_{<i}:=(x_1,\dots,x_{i-1}),
\]
而当 $i=1$ 时，$p(x_1\mid x_{<1})$ 直接理解为 $p(x_1)$。

这一恒等式是精确的。它本身不包含任何近似，也不施加任何建模假设。只要选定变量顺序，任意联合分布都可以用这种方式分解。

真正的建模步骤在于对条件密度进行参数化：
\[
p_\theta(x)
=
\prod_{i=1}^d p_\theta(x_i\mid x_{<i}).
\]
如果每个条件密度都已归一化，那么它们的乘积就定义了一个正确归一化的联合密度。这正是自回归模型按构造即可获得精确似然的原因。

\begin{tcolorbox}[
  colback=blue!4,
  colframe=blue!50!black,
  title={为什么自回归乘积是有效的概率密度},
  boxrule=0.6pt,
  arc=2mm,
  left=2mm,
  right=2mm,
  top=1mm,
  bottom=1mm
]
假设每个条件密度 $p_\theta(x_i\mid x_{<i})$ 都归一化：
\[
\int p_\theta(x_i\mid x_{<i})\,dx_i=1.
\]
那么联合密度
\[
p_\theta(x_1,\dots,x_d)
=
\prod_{i=1}^d p_\theta(x_i\mid x_{<i})
\]
也同样归一化。从 $x_d$ 开始向前积分，首先有
\[
\int p_\theta(x_d\mid x_{<d})\,dx_d=1,
\]
接着有
\[
\int p_\theta(x_{d-1}\mid x_{<d-1})\,dx_{d-1}=1,
\]
依此类推，直到只剩下 $p_\theta(x_1)$，而它的积分同样为一。

因此，自回归乘积不仅是一种方便的分解形式；只要每个条件项本身是有效的，它就定义了一个有效的概率密度。
\end{tcolorbox}

\subsection{条件参数化与生成}

一旦规定如何表示每个条件密度，链式法则分解就成为一个可学习的生成模型。一般写作
\[
p_\theta(x_i\mid x_{<i}),
\qquad i=1,\dots,d,
\]
其中 $\theta$ 表示模型参数。

需要强调的是，这些条件分布通常并不是由 $d$ 个互不相关的模型分别表示。相反，通常使用一个统一的参数化函数——往往是神经网络——在遵守自回归顺序的同时，输出所有条件分布的参数。整个分解过程中复用的是同一组模型参数；从一个条件分布到下一个条件分布，变化的是条件上下文。

对于连续变量，一个常见而简单的选择是高斯条件分布：
\[
p_\theta(x_i\mid x_{<i})
=
\mathcal{N}\!\left(
x_i;
\mu_i(x_{<i}),
\sigma_i^2(x_{<i})
\right),
\]
其中均值和方差都是先前变量的函数。更灵活的条件分布还可以采用高斯混合、离散化 Logistic 混合或其他富有表达能力的一维分布。

对于离散变量，模型可以输出类别分布的概率：
\[
p_\theta(x_i\mid x_{<i})
=
\operatorname{Categorical}\!\left(\pi_i(x_{<i})\right).
\]
无论哪种情形，模型预测的都是完整的条件分布，而不仅是一个点估计。

\begin{tcolorbox}[
  colback=blue!4,
  colframe=blue!50!black,
  title={三变量自回归示例},
  boxrule=0.6pt,
  arc=2mm,
  left=2mm,
  right=2mm,
  top=1mm,
  bottom=1mm
]
对于 $x=(x_1,x_2,x_3)$，链式法则给出
\[
p_\theta(x_1,x_2,x_3)
=
p_\theta(x_1)\,
p_\theta(x_2\mid x_1)\,
p_\theta(x_3\mid x_1,x_2).
\]

若采用高斯条件分布，则每个条件分布都可以由一个均值和一个方差表示。共享的参数生成函数 $f_\theta$ 可以输出这些参数：
\[
(\mu_1,\sigma_1^2)=f_\theta(\varnothing),
\]
\[
(\mu_2,\sigma_2^2)=f_\theta(x_1),
\]
\[
(\mu_3,\sigma_3^2)=f_\theta(x_1,x_2).
\]

这里 $f_\theta(\varnothing)$ 表示第一个条件分布没有任何先前变量可以作为条件，因此它输出 $p_\theta(x_1)$ 的无条件参数。随后，$f_\theta(x_1)$ 使用已知的 $x_1$ 生成 $p_\theta(x_2\mid x_1)$ 的参数，而 $f_\theta(x_1,x_2)$ 则利用前两个变量生成 $p_\theta(x_3\mid x_1,x_2)$ 的参数。

例如，
\[
p_\theta(x_2\mid x_1)
=
\mathcal{N}\!\left(x_2;\mu_2,\sigma_2^2\right),
\qquad
(\mu_2,\sigma_2^2)=f_\theta(x_1),
\]
同理，
\[
p_\theta(x_3\mid x_1,x_2)
=
\mathcal{N}\!\left(x_3;\mu_3,\sigma_3^2\right),
\qquad
(\mu_3,\sigma_3^2)=f_\theta(x_1,x_2).
\]
\end{tcolorbox}

整个过程中复用的是同一组模型参数 $\theta$。从一个条件分布到下一个条件分布，改变的不是模型本身，而是当前步骤可用的条件信息量。这可以消除一个常见误解：自回归分解虽然列出了一系列条件分布，却并不意味着我们为每个因子训练一个彼此无关的独立模型。在现代自回归模型中，这些条件分布通常由一个具有受限信息流的共享架构统一表示。

同一个分解也自然给出了直接的采样过程。给定
\[
p_\theta(x)
=
\prod_{i=1}^d p_\theta(x_i\mid x_{<i}),
\]
可以按顺序生成样本：
\[
x_1\sim p_\theta(x_1),
\]
然后
\[
x_2\sim p_\theta(x_2\mid x_1),
\]
再然后
\[
x_3\sim p_\theta(x_3\mid x_1,x_2),
\]
依此类推，直到
\[
x_d\sim p_\theta(x_d\mid x_1,\dots,x_{d-1}).
\]
所有分量生成后，所得向量
\[
x=(x_1,\dots,x_d)
\]
就是来自该自回归模型的一个样本。

应把这一过程理解为完整的生成机制。模型并不是先采样一个潜变量再对其解码，也不是运行随机去噪过程，而是直接逐分量构造数据点，并使用已经生成的分量作为后续分量的上下文。

\subsection{并行训练、顺序采样与掩码}

自回归模型具有一种重要的计算不对称性：训练期间的似然计算通常可以并行化，而采样在本质上是顺序的。

给定一个观测数据点
\[
x=(x_1,\dots,x_d),
\]
其对数似然为
\[
\log p_\theta(x)
=
\sum_{i=1}^d
\log p_\theta(x_i\mid x_{<i}).
\]
训练时，$x$ 的全部分量都已知。因此，每个条件上下文 $x_{<i}$ 都可以直接从观测数据中读取，无需采样。模型使用真实的先前分量来计算每个条件项。

这就是自回归似然训练精确且稳定的原因。模型不需要先生成 $x_1$，再用生成出的值来计算 $x_2$，如此继续；相反，它直接针对观测数据计算所有条件概率：
\[
\log p_\theta(x_1),
\quad
\log p_\theta(x_2\mid x_1),
\quad
\dots,
\quad
\log p_\theta(x_d\mid x_{<d}).
\]
具体取决于架构，其中许多乃至全部条件项都可以并行计算。

采样则不同。生成新样本时，未来分量尚未产生，因此是未知的。于是必须按顺序进行：
\[
x_1\sim p_\theta(x_1),
\]
\[
x_2\sim p_\theta(x_2\mid x_1),
\]
\[
\cdots
\]
\[
x_d\sim p_\theta(x_d\mid x_{<d}).
\]
每个新采样得到的分量都会成为下一步条件上下文的一部分，因此采样从构造上就是顺序执行的。

为了在神经网络中实现自回归分解，还必须保证第 $i$ 个条件分布的参数只能依赖 $x_{<i}$，不能依赖 $x_i$ 本身或未来变量。这一限制通常通过\emph{掩码（masking）}来强制实现。掩码会阻断模型中的非法连接，确保信息只能从较早变量流向较晚变量。

因此，即使训练时网络一次处理整个向量 $x$，针对 $x_i$ 输出的参数也不能依赖 $x_i,x_{i+1},\dots,x_d$。计算可以并行，但依赖结构仍然保持有序。

抽象地说，掩码施加了一个下三角依赖模式。如果定义依赖矩阵 $D$，其中 $D_{ij}$ 表示 $x_i$ 的条件分布是否允许依赖 $x_j$，那么自回归顺序要求
\[
D_{ij}=0
\qquad
\text{只要}
\qquad
j\geq i.
\]
例如，
\[
D
=
\begin{pmatrix}
0      & 0      & 0      & \cdots & 0 \\
\ast   & 0      & 0      & \cdots & 0 \\
\ast   & \ast   & 0      & \cdots & 0 \\
\vdots & \vdots & \vdots & \ddots & 0 \\
\ast   & \ast   & \ast   & \cdots & 0
\end{pmatrix},
\]
其中 $\ast$ 表示允许存在的依赖关系。

这种三角依赖是模型在实践中保持可处理性的关键原因之一。它避免了循环依赖：$x_i$ 的分布可以依赖更早的分量，但更早分量不会同时反过来依赖 $x_i$。因此，高维联合似然可以分解为条件对数似然之和。

从这个意义上说，自回归模型中的三角依赖，与基于耦合的归一化流中的三角 Jacobian 起着相似作用。两种情况下，依赖关系都沿单一方向组织，使高维概率计算可以分解为简单项。

掌握这一原则之后，许多著名架构都可以在结构层面得到理解。MADE 通过掩码在前馈网络中强制自回归结构~\cite{germain2015made}；PixelRNN 和 PixelCNN 将自回归分解用于图像~\cite{oord2016pixel}；WaveNet 则利用因果卷积生成原始音频~\cite{oord2016wavenet}。

因此，自回归模型提供了第二条实现精确密度建模的重要路径。归一化流依靠可逆变换和可处理的 Jacobian 行列式获得精确似然；自回归模型则依靠有序条件分解获得精确似然。二者都表明：当高维概率模型中的依赖关系经过精心组织时，其计算仍然可以保持可处理。


\section*{总结与参考文献}
\addcontentsline{toc}{section}{总结与参考文献}

本章介绍了精确密度建模这一与前文潜变量模型和扩散模型互补的范式。在这一范式中，似然可以直接求值，不需要变分下界、随机加噪过程或学习得到的得分场。

我们首先研究了归一化流。其核心思想是使用可微、可逆的映射，把简单的基础分布变换为复杂的数据分布。变量变换公式负责追踪局部体积膨胀或压缩如何改变概率密度，由此引出了流层的主要设计要求：可逆性、维度保持以及可处理的 Jacobian 行列式。

随后我们考察了若干代表性流构造。平面流展示了简单的非线性秩一扰动如何在保持 Jacobian 行列式可处理的同时引入形变。仿射耦合层则提供了更实用的机制：每次只变换输入的一部分，并以其余部分作为条件，从而同时获得闭式逆映射和高效的对数行列式计算。将这些层与坐标混合操作共同堆叠，就可以让简单的局部变换形成具有精确似然计算和精确模型采样能力的高表达生成模型。

本章后半部分介绍了自回归模型，将其作为另一条获得精确似然的路径。自回归模型不通过可逆映射变换密度，而是利用概率链式法则把联合密度分解为条件分布的乘积。这样无需潜变量、变分下界或 Jacobian 行列式，也能获得精确似然。代价是必须规定一个顺序，而且采样通常是顺序执行的。我们还看到，掩码如何在神经网络架构中强制所需的有序信息流，使训练时能够高效计算似然项，同时保持自回归分解成立。

贯穿这两类模型的中心观点是：精确似然之所以成为可能，依赖于结构约束。在归一化流中，这些约束体现为可逆性与可处理的 Jacobian 行列式；在自回归模型中，则体现为有序的条件依赖。两种方法都说明，只要对变换或依赖关系进行精心组织，高维概率模型仍可以保持计算可行。

\bigskip

Rezende 和 Mohamed~\cite{rezende2015variational} 将归一化流作为通过一系列可逆变换构造灵活分布的一般框架加以提出。重要的实用流架构包括 Dinh 等人提出的 NICE 和 RealNVP~\cite{dinh2014nice,dinh2017density}，它们发展了具有可处理 Jacobian 的耦合式变换；以及 Kingma 与 Dhariwal 提出的 Glow~\cite{kingma2018glow}，后者通过可逆 $1\times1$ 卷积和可扩展流设计扩展了这些思想。Papamakarios 等人~\cite{papamakarios2021normalizing} 则给出了关于归一化流及其在密度估计中作用的更广泛综述。

自回归神经密度估计同样在现代生成建模中发挥了重要作用。Germain 等人的 MADE~\cite{germain2015made} 展示了如何通过掩码在前馈神经网络中强制自回归结构。van den Oord 等人的 PixelRNN 和 PixelCNN~\cite{oord2016pixel} 发展了自回归图像模型，而 WaveNet~\cite{oord2016wavenet} 则把因果自回归建模用于原始音频。综合来看，这些工作体现了基于精确似然的建模方法持续的重要性——无论这种精确性来自可逆变换，还是来自有序条件分解。
